\section{Ethics, safety, and resources}

\subsection{Ethical considerations}
This work studies mechanistic aspects of research workflow design rather than
reporting a deployment in a human-facing domain. The frozen local evaluation
uses constructed cases and does not require human-subject records or personally
identifiable information. Any later external-provider or retrieval execution
must be documented separately from this local evidence rather than inferred from
the protocol design.

\subsection{Safety considerations}
Authority promotion and formulation reconstruction are separate concerns, and
only the second is studied here. The mechanisms described can expose a possible
method failure but cannot authorize $M_t$ to rewrite itself. The reframe gate is
mechanically bounded: only formulation and search-universe coordinates may
trigger a reframe, and the evidence-only and execution-only negative controls
exercise the corresponding non-reframe behaviour in the frozen local suite.

\subsection{Resource usage and environmental impact}
The repository does not yet contain an admissible clean-environment resource
ledger for the full campaign. Numerical model-token, runtime, monetary-cost and
carbon-impact quantities therefore remain undetermined in this revision. A
future result-bearing run must retain actual resource usage under its frozen
provider and environment bindings before those quantities are reported.

\subsection{Dual-use considerations}
The epistemic-state machinery described here is general-purpose and could be
applied outside scientific research. This paper evaluates scientific-workflow
mechanics and does not present an empirical evaluation of military or other
high-risk deployment applications. The authority boundaries above are part of
the implemented design, but they should not be read as evidence of safety in
domains that were not evaluated.
